problem:  
Let \(G = SU(3)\), and let \(P\subset G\) be the standard \(U(2)\subset SU(3)\).  
Consider the one-parameter family of left-invariant metrics \(g_\lambda\) on \(SU(3)\) defined by
\[
g_\lambda(X,Y) =
\begin{cases}
\langle X,Y\rangle_0, & X,Y\in \mathfrak{p} = \mathrm{Lie}(P),\\
\lambda\,\langle X,Y\rangle_0, & X,Y\in \mathfrak{p}^{\perp},
\end{cases}
\]
where \(\langle\cdot,\cdot\rangle_0\) is the bi-invariant metric arising from the negative Killing form and \(0<\lambda<1\).  
For integer triples \(p=(p_1,p_2,p_3)\), \(q=(q_1,q_2,q_3)\) with \(\sum p_i = \sum q_i\) and free circle action
\[
(z,g)\mapsto \operatorname{diag}(z^{p_1},z^{p_2},z^{p_3})\, g\, \operatorname{diag}(z^{-q_1},z^{-q_2},z^{-q_3}),
\]
let \(E_{p,q}=S^1_{p,q}\backslash SU(3)\) with the quotient metric \(g^{(\lambda)}_{p,q}\) induced from \(g_\lambda\).  
For many parameter pairs \((p,q)\) satisfying Eschenburg’s positivity inequalities, \(g^{(\lambda)}_{p,q}\) has positive sectional curvature.

Now fix some particular \(\lambda=\lambda_*\) within a positive-curvature range, and ask:

> **Rigidity problem.**  
> Suppose two Eschenburg spaces \(E_{p,q}\) and \(E_{p',q'}\) are diffeomorphic as smooth manifolds (as verified, for example, by equality of all Kreck–Stolz invariants) and each admits positive curvature for the same fixed metric parameter \(\lambda_*\).  
> Are the corresponding canonical quotient metrics \(g^{(\lambda_*)}_{p,q}\) and \(g^{(\lambda_*)}_{p',q'}\) necessarily isometric under some automorphism of \(SU(3)\) together with reparametrization of the \(S^1\) action, or can distinct parameter pairs yield diffeomorphic manifolds that carry genuinely different (\emph{non-isometric}) canonical positively curved metrics?

Why is it a "good" problem:  
This formulation fully specifies both the metric class and the equivalence relations, eliminating earlier ambiguities. It asks a concrete rigidity question within a tightly controlled, explicit family of examples where tools from Lie groups, curvature computation, and diffeomorphism classification can all be brought to bear.  

The problem sits precisely at the intersection of Riemannian geometry (uniqueness and moduli of positively curved metrics), topology (classification via Kreck–Stolz invariants), and symmetry analysis (automorphisms of \(SU(3)\)).  
A positive resolution would reveal deep rigidity linking geometric curvature data to topological type; a negative one would expose genuine moduli of positive curvature metrics—something scarcely known to exist.  
Because the question is cleanly stated, technically well-posed, and probes a boundary phenomenon in the theory of positive curvature, it exemplifies a “good” mathematical problem: simple to state, profound in implication, and bridging geometry, topology, and Lie theory.